%-------------------------
% Resume in Latex
% Author : Harshibar
% Based off of: https://github.com/jakeryang/resume
% License : MIT
%------------------------

\documentclass[letterpaper,11pt]{article}

\usepackage{latexsym}
\usepackage[empty]{fullpage}
\usepackage{titlesec}
\usepackage{marvosym}
\usepackage[usenames,dvipsnames]{color}
\usepackage{verbatim}
\usepackage{enumitem}
\usepackage[hidelinks]{hyperref}
\usepackage{fancyhdr}

\usepackage[T2A]{fontenc}
\usepackage[utf8]{inputenc}
\usepackage[english,russian]{babel}


\usepackage{tabularx}
% only for pdflatex
% \input{glyphtounicode}

% fontawesome
% \usepackage{fontawesome5}
\usepackage{fontawesome}

% fixed width
% \usepackage[scale=0.90,lf]{FiraMono}

% light-grey
\definecolor{light-grey}{gray}{0.83}
\definecolor{dark-grey}{gray}{0.3}
\definecolor{text-grey}{gray}{.08}

\DeclareRobustCommand{\ebseries}{\fontseries{eb}\selectfont}
\DeclareTextFontCommand{\texteb}{\ebseries}

% custom underilne
\usepackage{contour}
\usepackage[normalem]{ulem}
\renewcommand{\ULdepth}{1.8pt}
\contourlength{0.8pt}
\newcommand{\myuline}[1]{%
  \uline{\phantom{#1}}%
  \llap{\contour{white}{#1}}%
}

% custom font: helvetica-style
% \usepackage{tgheros}
% \renewcommand*\familydefault{\sfdefault} 
% %% Only if the base font of the document is to be sans serif
% \usepackage[T1]{fontenc}
\usepackage{fontspec}
% \setmainfont{Arial}
% \setsansfont{Arial}
% \setmainfont{Liberation Sans}
% \setsansfont{Liberation Sans}
\setmainfont{FreeSans}
\setsansfont{FreeSans}
\renewcommand{\familydefault}{\sfdefault}


\pagestyle{fancy}
\fancyhf{} % clear all header and footer fields
\fancyfoot{}
\renewcommand{\headrulewidth}{0pt}
\renewcommand{\footrulewidth}{0pt}

% Adjust margins
\addtolength{\oddsidemargin}{-0.5in}
\addtolength{\evensidemargin}{0in}
\addtolength{\textwidth}{1in}
\addtolength{\topmargin}{-.5in}
\addtolength{\textheight}{1.0in}

\urlstyle{same}

\raggedbottom
\raggedright
\setlength{\tabcolsep}{0in}

% Sections formatting - serif
% \titleformat{\section}{
%   \vspace{2pt} \scshape \raggedright\large % header section
% }{}{0em}{}[\color{black} \titlerule \vspace{-5pt}]

% TODO EBSERIES
% sans serif sections
\titleformat {\section}{
    \bfseries \vspace{2pt} \raggedright \large % header section
}{}{0em}{}[\color{light-grey} {\titlerule[2pt]} \vspace{-4pt}]

% only for pdflatex
% Ensure that generate pdf is machine readable/ATS parsable
% \pdfgentounicode=1

%-------------------------
% Custom commands
\newcommand{\resumeItem}[1]{
  \item\small{
    {#1 \vspace{-1pt}}
  }
}

\newcommand{\resumeSubheading}[4]{
  \vspace{-1pt}\item
    \begin{tabular*}{\textwidth}[t]{l@{\extracolsep{\fill}}r}
      \textbf{#1} & {\color{dark-grey}\small #2}\vspace{1pt}\\ % top row of resume entry
      \textit{#3} & {\color{dark-grey} \small #4}\\ % second row of resume entry
    \end{tabular*}\vspace{-4pt}
}

\newcommand{\resumeSubSubheading}[2]{
    \item
    \begin{tabular*}{\textwidth}{l@{\extracolsep{\fill}}r}
      \textit{\small#1} & \textit{\small #2} \\
    \end{tabular*}\vspace{-7pt}
}

\newcommand{\resumeProjectHeading}[2]{
    \item
    \begin{tabular*}{\textwidth}{l@{\extracolsep{\fill}}r}
      #1 & {\color{dark-grey}} \\
    \end{tabular*}\vspace{-4pt}
}

\newcommand{\resumeSubItem}[1]{\resumeItem{#1}\vspace{-4pt}}

\renewcommand\labelitemii{$\vcenter{\hbox{\tiny$\bullet$}}$}

% CHANGED default leftmargin  0.15 in
\newcommand{\resumeSubHeadingListStart}{\begin{itemize}[leftmargin=0in, label={}]}
\newcommand{\resumeSubHeadingListEnd}{\end{itemize}}
\newcommand{\resumeItemListStart}{\begin{itemize}}
\newcommand{\resumeItemListEnd}{\end{itemize}\vspace{0pt}}

\color{text-grey}

%-------------------------------------------
%%%%%%  RESUME STARTS HERE  %%%%%%%%%%%%%%%%%%%%%%%%%%%%

\begin{document}

%----------HEADING----------
\begin{center}
    \textbf{\Huge Милевская Виктория Андреена } \\ \vspace{5pt}
    \small \texttt{+7 (916) 744-21-25} \hspace{1pt} \textbar
    \hspace{1pt} \texttt{vicktoriamilewskaya@yandex.ru}
    \\ \vspace{-3pt}
\end{center}

%-----------EXPERIENCE-----------
\section{Опыт}
  \resumeSubHeadingListStart

    \resumeSubheading
      {Яндекс}{Август 2022 — сейчас}
      {Разработчик интерфейсов}{}
      \resumeItemListStart
        \resumeItem{Конструктор «Нейролендинги»: разработала полноценный продукт конуструктор для генерации лендингов с редактором блоков, модулем управления товарами и системой кастомизации. Спроектировала архитектуру клиентской части, сделала интерфейс, логику генерации и API}
        \resumeItem{Разработала сценарий рекламы в Поиске специалистов и услуг: реализовала карточки в поиске, а так же сценарий в рекламном кабинете.}
        \resumeItem{Объединила дежурства 2х отделов, привела дежурство к единому флоу, реализовала техдолг и инструменты для поддержки сервиса, что сократило количество тикетов, доходящих до разработки сервиса, в 4 раза.}
         \resumeItem{Генеративные материалы в рекламном кабинете: реализовала внедрение саджеста генеративных картинок в рекламный кабинет, переписала полностью управление картинками в рекламе. Использование картинок, сгенерированные моделью, привело к тренду использования генеративных материалов ко всем рекламным материалам. Добавила саджест целевых действий, снизив долю кампаний без конверсий с 69 процентов до 67. }
        \resumeItem{Инфраструктура сервиса: провела масштабное обновление балансировщиков для микросервисов приложения, переведя их на новую архитектуру, что исключило падения во время инцидентов в ДЦ компании.}
      \resumeItemListEnd

    \resumeSubheading
      {Сбер}{Май 2021 — Август 2022}
      {Инженер-программист}{}
      \resumeItemListStart
        \resumeItem{Полный цикл разработки клиент-серверных приложений: проектирование архитектуры БД, разработка REST API и создание адаптивных пользовательских интерфейсов.}
        \resumeItem{Разработка высоконагруженной системы мониторинга СЦ: реализовала интерфейсы для прогнозирования инцидентов, графы связей и интеллектуальные дашборды (включая кастомизацию плагинов для Grafana)}
        \resumeItem{Разработка внутренней антифрод-системы: реализовала  интерфейсы для прогнозирования инцидентов, графы связей и интеллектуальные дашборды (включая кастомизацию плагинов для Grafana)}
    \resumeItemListEnd

    \resumeSubheading
      {Comindware}{Февраль 2020 — Май 2021}
      {Инженер-программист}{}
      \resumeItemListStart
        \resumeItem{Разработка мессенджера в составе Enterprise-платформы BPM/CRM-системы. Разработала архитектуру клиентского приложения: спроектировала модули для работы с сообщениями (отправка, получение, статусы, история), навигационные компоненты и контекстную информационную панель. Функционал работы с сообщениями такие как реплаи и форварды. Выполнила интеграцию React-компонентов в существующую Backbone-инфраструктуру, обеспечив бесшовную работу гибридного интерфейса.}
        \resumeItem{Реализовала оптимизацию продукта для скринридеров.}
        \resumeItem{Выполнила редизайн админской части приложения.}
      \resumeItemListEnd

  \resumeSubHeadingListEnd


%-----------EDUCATION-----------
\section {Образование}
  \resumeSubHeadingListStart
    \resumeSubheading
      {Национальный исследовательский университет «Высшая школа экономики»}{2025 -- сейчас}
      {бакалавриат прикладной математики и информатики}{Москва}
  \resumeSubHeadingListEnd


%
%-----------PROGRAMMING SKILLS-----------
\section{Навыки}
 \begin{itemize}[leftmargin=0in, label={}]
    \small{\item{
    \textbf{Языки} {: Русский, Английский}\vspace{2pt} \\
     \textbf{Python, Node.js, JavaScript, TypeScript, C++, HTML5, CSS/SCSSS, SQL (PostgreSQL, MySQL)}\vspace{2pt} \\
      Git, ES6, React, Redux, MobX, Jest, Ajax, GraphQL, WebSocket, Django, Flask, FastAPI, SQLAlchemy, Jinja\vspace{2pt} \\
      Постановка задач, декомпозиция задач, сode review, умение работать в команде, умение разбираться в чужом коде
    }}
 \end{itemize}


%-------------------------------------------
\end{document}